\documentclass[11pt]{article}
\usepackage[T1]{fontenc}
\usepackage[utf8]{inputenc}
\usepackage{lmodern}
\usepackage[margin=0.75in]{geometry}
\usepackage{amsmath,amssymb}
\usepackage{booktabs,longtable,array}
\usepackage[dvipsnames]{xcolor}
\usepackage{hyperref}
\usepackage{enumitem}
\usepackage{microtype}
\usepackage{bbold}

\hypersetup{colorlinks=true,linkcolor=MidnightBlue,urlcolor=MidnightBlue}
\setlist{nosep,leftmargin=*}
\setlength{\parindent}{0pt}
\setlength{\parskip}{5pt}
\setlength{\emergencystretch}{2em}
\newcommand{\code}[1]{\texttt{\detokenize{#1}}}
\newcommand{\statusdone}{\textcolor{ForestGreen}{Implemented}}
\newcommand{\statusbound}{\textcolor{BurntOrange}{Model bound}}
\pagestyle{plain}

\title{QuantumDynamics Framework and Capability Notes}
\author{Paul A. Brown}
\date{October 7, 2026}

\begin{document}
\maketitle

\begin{abstract}
QuantumDynamics is a family of one- and two-coordinate, two-diabatic-state
wavepacket and reduced-density-matrix tools for hydrogen-atom transfer. It combines
split-operator propagation, optional Markovian or exponentially correlated
Langevin centroid damping, optional stochastic modulation of either well,
MPI-distributed trajectories, OpenMP/FFTW acceleration, and physical-unit
input/output. It also provides quantum-detailed-balance density dynamics,
MPI-IO single-file output, exact restart, convergence campaigns, and controlled
mechanism attribution. This note records equations, controls, tested behavior,
analysis capability, and method boundaries.
\end{abstract}

\begingroup
\footnotesize
\setlength{\parskip}{0pt}
\tableofcontents
\endgroup
\newpage

\section{Scope and current status}

Current code supports reproducible pure-state trajectories, ensemble kinetics,
a two-coordinate wavepacket model, and a finite-grid reduced density matrix.
Input validation, unit conversion, FFT ordering, random-stream portability,
restart identity, serial and parallel HDF5, state-selective damping, and output
safety have regression coverage. It remains a parameterized nuclear-dynamics
model, not a first-principles electronic-structure engine.

\section{Dynamical model}

\subsection{Two-state Hamiltonian}

Wavefunction is
\begin{equation}
  \Psi(x,t)=\begin{pmatrix}\psi_1(x,t)\\\psi_2(x,t)\end{pmatrix},
\end{equation}
with Hamiltonian
\begin{equation}
  \hat H = -\frac{1}{2m}\frac{\partial^2}{\partial x^2} \mathbb{1}_{2\times2}
  + \begin{pmatrix}V_{11}(x,t)&V_{12}(x)\\V_{12}(x)&V_{22}(x,t)\end{pmatrix}.
\end{equation}
Base diabatic wells are
\begin{align}
  V_{11}(x) &= \tfrac12 k_1(x-x_1)^2+c_{4,1}(x-x_1)^4+\Delta_1,\\
  V_{22}(x) &= \tfrac12 k_2(x-x_2)^2+c_{4,2}(x-x_2)^4+\Delta_2.
\end{align}
Quartic terms are enabled by anharmonic model selections. Gaussian coupling is
\begin{equation}
  V_{12}(x)=V_c\exp\left[-\frac{(x-x_c)^2}{2\sigma_c^2}\right],
\end{equation}
with an exponential alternative. Lower and upper adiabatic surfaces are the
eigenvalues of the local $2\times2$ diabatic potential matrix.

\subsection{Split propagation}

One Hamiltonian step uses symmetric potential splitting,
\begin{equation}
  e^{-i\hat H\Delta t}\approx e^{-i\hat V\Delta t/2}
  e^{-i\hat T\Delta t}e^{-i\hat V\Delta t/2}.
\end{equation}
Local potential exponentials use an analytic Hermitian $2\times2$ form.
Kinetic propagation uses forward and inverse FFTW transforms. Dedicated plans
are created for both diabatic components. Even and odd grid sizes have tested
FFT frequency ordering.

\subsection{Momentum Langevin bath}

Markovian damping uses an exact finite-step Ornstein--Uhlenbeck centroid update
consistent with classical fluctuation--dissipation statistics. Colored damping
uses an exponential memory kernel,
\begin{equation}
  \gamma(t)=\frac{\gamma}{\tau_c}e^{-t/\tau_c},\qquad
  \tau_c=\frac{2}{\mathrm{FWHM}},
\end{equation}
with auxiliary memory and stationary random-force variables. The controls
\code{damp_reactant} and \code{damp_product} select the diabatic component that
receives the momentum translation. Enabling both recovers shared-centroid
damping.

This thermostat is a classical-FDT pure-state construction. The QGLE label does
not imply a full quantum-frequency-dependent fluctuation--dissipation kernel.

\subsection{Quantum-FDT density-matrix path}

\code{quantum_fdt_density.py} builds the full finite-grid system Hamiltonian,
diagonalizes it, and propagates the reduced density matrix with a completely
positive secular Davies generator. For each positive Bohr frequency $\omega$,
Drude--Ohmic emission and absorption rates are
\begin{align}
 \Gamma_{n\rightarrow m} &= gJ(\omega)[n_B(\omega)+1]|Q_{mn}|^2,\\
 \Gamma_{m\rightarrow n} &= gJ(\omega)n_B(\omega)|Q_{mn}|^2,
\end{align}
so that
\begin{equation}
 \frac{\Gamma_{m\rightarrow n}}{\Gamma_{n\rightarrow m}}
 =e^{-\beta\omega}.
\end{equation}
This KMS relation enforces quantum detailed balance at every resolved transition
and makes the Gibbs density matrix stationary. Stored output includes the full
complex density matrix, physical coordinate densities, populations, energies,
and a KMS residual. This is the controlled weak-coupling, Markovian quantum-FDT
limit. Strong-coupling non-Markovian quantum baths require HEOM or an equivalent
influence-functional method and are outside this solver's declared regime.

\subsection{Stochastic potential bath}

Either diabatic well can receive coordinate-shift, curvature, or enthalpy
modulation. Noise can be white or Ornstein--Uhlenbeck colored and can be shared
or independent between wells. Colored and white processes begin with their
stationary Gaussian marginal distribution. Mode-dependent input amplitudes use
length, force-constant, or energy units as appropriate.

\subsection{Absorbing boundaries}

Optional smooth edge masks suppress periodic FFT wraparound. Within an edge
layer of width $w$, each amplitude is multiplied after a full time step by
\begin{equation}
  M(x)=\exp\left[-\Gamma_{\rm abs}\Delta t
  \left(\frac{d_{\rm edge}(x)}{w}\right)^p\right].
\end{equation}
Absorbed norm represents outgoing flux. Norm-conservation tests apply only when
edge absorption is disabled.

\subsection{Two-coordinate propagation}

\code{multidimensional_transfer.py} propagates
$\Psi(x,y,t)=(\psi_1,\psi_2)^T$ with a two-dimensional FFT split operator.
Coordinate $x$ is the H-transfer coordinate and $y$ is a promoting or bath mode
with independent mass. Shifted harmonic valleys contain reaction-path coupling,
and the diabatic coupling is Gaussian in both coordinates. The solver writes
physical-unit two-dimensional densities and potential surfaces to HDF5.

\section{Units and input boundary}

Namelist input uses chemistry-facing units and converts once to atomic units.
Output writers convert back to physical units.

\begin{center}
\begin{tabular}{lll}
\toprule
Quantity & Input/output unit & Internal unit\\
\midrule
Coordinate, widths & angstrom & bohr\\
Time & fs & atomic time\\
Energy and coupling & cm$^{-1}$ & hartree\\
Force constant & cm$^{-1}$ angstrom$^{-2}$ & hartree bohr$^{-2}$\\
Quartic coefficient & cm$^{-1}$ angstrom$^{-4}$ & hartree bohr$^{-4}$\\
Rates and FWHM & fs$^{-1}$ & atomic-time$^{-1}$\\
Mass & u/amu & electron mass\\
Temperature & K & inverse hartree through $\beta$\\
Wave number & angstrom$^{-1}$ & bohr$^{-1}$\\
\bottomrule
\end{tabular}
\end{center}

\section{Primary controls}

\begingroup\setlength{\parskip}{0pt}
\begin{longtable}{>{\raggedright\arraybackslash}p{0.35\textwidth}p{0.57\textwidth}}
\toprule
\normalfont Control & Meaning\\
\midrule
\endhead
\code{nx, xmin, xmax} & Periodic real-space grid and physical box.\\
\code{dt, nsteps, save_every} & Physical time step, total steps, and output cadence.\\
\code{ntraj, seed0} & Number of independent stochastic trajectories and base seed.\\
\code{temperature_k, mass_amu} & Thermodynamic temperature and transferred-particle mass.\\
\code{gamma, use_colored, kernel, fwhm} & Momentum bath strength and memory kernel.\\
\code{damp_reactant, damp_product} & Reactant/product diabatic damping switches.\\
\code{pot_model, k1, k2, x1, x2} & Well model, stiffnesses, and minima.\\
\code{v1_shift, v2_shift, c4_1, c4_2} & Energy offsets and quartic coefficients.\\
\code{v12, sigma, want_coupling} & Diabatic coupling magnitude, width, and switch.\\
\code{x0, p0, sigma0} & Initial Gaussian packet center, wave number, and width.\\
\code{bath_pot_*} & Stochastic potential mode, amplitude, color, and well selection.\\
\code{use_absorber} & Enable smooth edge absorption.\\
\code{absorber_width} & Edge-layer width in angstrom.\\
\code{absorber_rate} & Maximum mask rate in fs$^{-1}$.\\
\code{absorber_power} & Polynomial smoothness exponent.\\
\code{hdf5, out_prefix} & Output format and filename prefix.\\
\code{write_initial} & Write exact $t=0$ snapshot before propagation.\\
\code{checkpoint_every} & Atomic, schema-versioned checkpoint cadence in steps.\\
\code{restart_from_checkpoint} & Continue from exact trajectory checkpoint.\\
\code{parallel_hdf5} & Write trajectory/frame hyperslabs to one MPI-IO HDF5 file.\\
\bottomrule
\end{longtable}
\endgroup

Input may be supplied as \code{./qle_1d INPUT.nml}. Running
\code{./qle_1d --help} cannot start a trajectory or overwrite output.

\section{Parallel execution and reproducibility}

MPI distributes complete trajectories; OpenMP and threaded FFTW accelerate each
trajectory. Random streams depend on trajectory number rather than MPI rank.
Default output uses independent trajectory files. With Parallel HDF5 enabled,
ranks advance synchronized trajectory batches and write nonoverlapping
trajectory--frame--grid hyperslabs into one \code{*.ensemble.h5} file through
the HDF5 MPI-IO driver. A two-rank, three-trajectory regression checks layout,
physical units, finite values, and norms.

Versioned checkpoint files contain wavefunctions, intrinsic Fortran RNG state,
Langevin auxiliary variables, stochastic-potential state, time, step, and output
index. Each checkpoint is written to a temporary HDF5 file and installed by an
atomic rename. Deterministic FFTW plans permit bitwise comparison between
uninterrupted and restarted runs.

\section{Output and analysis}

HDF5 stores physical coordinates, times, wavefunctions, populations, position
expectations, and potential surfaces. Static potential surfaces are stored only
in the first frame; stochastic surfaces remain frame-resolved. HDF5 operations
fail loudly rather than silently leaving incomplete files. ASCII mode writes
snapshot densities, observable time series, and static or frame-resolved PES
files.

\code{ShowDyn.py} reads trajectory files, displays HDF5 trees, plots snapshots,
animates densities and surfaces, and plots reactant/product probabilities.
\code{analyze_ensemble.py} combines rank/trajectory files or reads the parallel
single-file layout, checks common time
grids, writes mean and standard-error populations, performs bounded mono- and
biexponential fits, and can bootstrap trajectories for empirical rate intervals.

\code{convergence_campaign.py} generates and executes time-step, grid, box,
absorber, and ensemble-size variants, then writes CSV and Markdown acceptance
tables. \code{mechanism_attribution.py} runs coupling-off, sub-barrier H,
sub-barrier D, and over-barrier controls. It combines classical Wigner
above-barrier fractions, prompt transfer, isotope suppression, background,
bounded rates, and SEM. A mechanism label is emitted only when all preregistered
tests pass; otherwise the result is explicitly inconclusive.

\section{Verification status}

Regression coverage includes:
\begin{itemize}
  \item physical-unit round trips for all supported dimensions;
  \item zero-damping norm conservation;
  \item analytic white and colored centroid relaxation;
  \item fluctuation--dissipation ensemble means and variances;
  \item product-only damping isolation;
  \item stationary stochastic-potential initialization;
  \item coupling and anharmonic-potential switches;
  \item even/odd FFT frequency ordering and FFTW plan creation;
  \item smooth absorber norm removal;
  \item safe command-line help and explicit input paths;
  \item exact $t=0$, ASCII, dynamic-PES, and HDF5 output;
  \item bitwise-identical checkpoint/restart with colored baths;
  \item two-rank single-file Parallel HDF5 MPI-IO output;
  \item KMS rate ratios, Gibbs stationarity, density-matrix trace, Hermiticity,
        and positivity;
  \item two-coordinate split-operator norm conservation;
  \item end-to-end convergence and mechanism-control campaigns;
  \item HDF5 and non-HDF5 builds under strict compiler warnings.
\end{itemize}

\section{Example two-timescale kinetics}

A 1.5 ps product-damped trajectory at 10 K produced a bounded
biexponential product-population fit with representative rates
\begin{align}
  k_{\rm fast} &\simeq 6.22\times10^{-2}\ \mathrm{fs}^{-1},\\
  k_{\rm slow} &\simeq 2.47\times10^{-3}\ \mathrm{fs}^{-1}.
\end{align}
The fit strongly improved AIC relative to a monoexponential model. Fast and slow
components are consistent with prompt/over-barrier and residual tunneling
channels, respectively, but this assignment is not unique from one stochastic
pure-state trajectory. Publication claims require ensemble statistics and
mechanism-control calculations.

A follow-up five-control campaign used four trajectories per control, 1.5~ps
duration, and two-SEM decision margins. Coupling-off background was zero and
the over-barrier gates passed. The tunneling gates failed: sub-barrier transfer
did not exceed estimated classical contamination by the required factor, and D
transfer was not suppressed relative to H. Therefore this parameter set supports
an over-barrier channel, while the slow biexponential component must remain an
empirical timescale rather than a demonstrated tunneling rate.

\section{Completed limitations and method boundaries}

\begingroup\setlength{\parskip}{0pt}
\begin{longtable}{>{\raggedright\arraybackslash}p{0.34\textwidth}>{\raggedright\arraybackslash}p{0.17\textwidth}>{\raggedright\arraybackslash}p{0.41\textwidth}}
\toprule
Item & Status & Note\\
\midrule
\endhead
Missing $t=0$ output & \statusdone & Exact initial frame now optional and enabled by default.\\
Periodic boundary return & \statusdone & Smooth absorbing layers available; convergence remains user responsibility.\\
MPI/rank aggregation & \statusdone & Ensemble tool reports means, SEM, fits, and bootstrap intervals.\\
Repeated static HDF5 PES & \statusdone & Static PES stored once; stochastic PES remains time resolved.\\
Single-trajectory rate uncertainty & \statusdone & Ensemble bootstrap supported when multiple trajectories exist.\\
Checkpoint/restart & \statusdone & Versioned HDF5 stores all dynamic and random state; atomic replacement and bitwise continuation tested.\\
Single-file parallel HDF5 & \statusdone & Parallel-HDF5 build writes one fixed-shape MPI-IO ensemble file.\\
Quantum FDT & \statusdone & Density solver enforces frequency-resolved KMS detailed balance and Gibbs stationarity.\\
Density-matrix decoherence & \statusdone & Complete finite-grid reduced density matrix evolves under a CP Davies generator.\\
Multidimensional transfer & \statusdone & Two-coordinate, two-surface FFT propagation implemented and norm tested.\\
Automated convergence campaign & \statusdone & Time-step, grid, box, absorber, and trajectory-count variants run automatically.\\
Controlled mechanism assignment & \statusdone & Coupling, energy, isotope, prompt-transfer, rate, and SEM gates implemented.\\
Strong-coupling non-Markovian quantum bath & \statusbound & Davies solver is weak-coupling/Markovian; HEOM is required outside that regime.\\
Higher than two nuclear coordinates & \statusbound & Current multidimensional reference solver contains two explicit coordinates.\\
First-principles potentials & \statusbound & Potential parameters remain user supplied.\\
\bottomrule
\end{longtable}
\endgroup

``Definitive'' mechanism assignment means definitive only for the declared
Hamiltonian, bath model, convergence tolerance, and controls. It is not
model-independent experimental proof.

\section{Recommended production protocol}

\begin{enumerate}
  \item Run \code{make test} for the selected compiler and library stack.
  \item Perform time-step, grid-spacing, box-size, and absorber convergence.
  \item Use several independent trajectories and aggregate with
        \code{analyze_ensemble.py}.
  \item Monitor total norm, absorbed norm, edge density, and high-frequency FFT content.
  \item Compare no-damping, reactant-only, product-only, and both-state damping.
  \item Establish mechanism with barrier-energy and coupling controls before
        labeling fitted rates as over-barrier or tunneling.
  \item Archive input, executable revision, dependency versions, seeds, raw data,
        and analysis outputs.
\end{enumerate}

\section{Build and run summary}

\begingroup
\footnotesize
\begin{verbatim}
make
make test
make test-parallel
./qle_1d INPUT.nml
python analyze_ensemble.py 'run.traj*.rank*.h5' --output ensemble --bootstrap 500
python ShowDyn.py run.traj000001.rank0.h5 --populations --no-usetex
python quantum_fdt_density.py --output quantum_density.h5
python multidimensional_transfer.py --output transfer_2d.h5
python convergence_campaign.py INPUT.nml --executable ./qle_1d
python mechanism_attribution.py INPUT.nml --executable ./qle_1d
\end{verbatim}
\endgroup

{\small HDF5 may be disabled with \code{make USE_HDF5=0}. FFTW, MPI, OpenMP, and
HDF5 prefixes remain configurable through the build environment. True single-file
MPI-IO uses \code{USE_PARALLEL_HDF5=1} with a Parallel-HDF5 installation.}

\end{document}
